\chapter{El juego (Sprint 2)}
\label{cap:juego}

\section{Concepto}

La evolución del proyecto en el Sprint 2 transformó la base SaaS inicial en un videojuego web de tipo \textit{survivor} (género popularizado por \textit{Vampire Survivors} \cite{vampiresurvivors}), ambientado en la historia y la cultura del Atlético de España. El jugador controla a un deportista histórico del club en un escenario bidimensional de perspectiva isométrica (cámara 3/4) donde las armas atacan de forma automática. El objetivo consiste en sobrevivir a oleadas continuas de rivales y jefes emblemáticos, recolectando experiencia, subiendo de nivel y eligiendo mejoras tácticas para prolongar la resistencia en el tiempo.

El proyecto mantiene su condición de iniciativa de aficionado sin vinculación oficial con el Club Atlético de Madrid, evitando el uso de elementos protegidos por propiedad intelectual y reinterpretando la iconografía rojiblanca a través de gráficos de estilo \textit{pixel art} generados algorítmicamente.

\section{Cómo se juega (manual de usuario)}

La experiencia de juego se estructura en torno a una serie de pantallas y mecánicas accesibles desde el menú principal:
\begin{enumerate}
    \item \textbf{Identificación y Acceso:} Con una sesión activa, la ruta \texttt{\#/jugar} da acceso al menú principal del juego. En el primer acceso se solicita un mote único y la elección de una cara para el avatar.
    \item \textbf{Menú del Juego:} Permite acceder al Vestuario (selección de jugadores y equipaciones), Entrenamiento, Contratar, Logros, Combinaciones, Dirección (entrenadores), Ranking y Ajustes.
    \item \textbf{Controles:} El desplazamiento se realiza mediante las teclas WASD o flechas del teclado, o mediante un joystick táctil flotante en dispositivos móviles. La pulsación de \texttt{Espacio} o \texttt{E} activa la habilidad definitiva (\textit{ulti}) del entrenador, \texttt{Esc} o \texttt{P} abre la pausa, \texttt{F} alterna la pantalla completa y las teclas numéricas \texttt{1} a \texttt{3} seleccionan las mejoras al subir de nivel.
    \item \textbf{Dinámica de Partida:} Cada 30 segundos se genera una nueva oleada de adversarios y cada dos minutos aparece un jefe principal. Los rivales derrotados otorgan experiencia; al acumular la suficiente, el nivel asciende y se despliega un panel con tres mejoras aleatorias.
    \item \textbf{Ambientes Temporales:} La partida evoluciona a lo largo del tiempo a través de cinco ambientes diferenciados que modifican el escenario cada cinco minutos: pretemporada, Final de la Copa del Rey, Final de Liga, Final de la Champions y Final del Mundial.
    \item \textbf{Fin de Partida:} Al agotarse los puntos de vida del personaje, se muestra una pantalla detallada con estadísticas de daño por habilidad, duración, bajas y recompensas obtenidas para el club.
\end{enumerate}

\section{Sistemas del juego}

El equilibrio y la variedad del juego descansan sobre múltiples sistemas interconectados definidos en el motor y en los ficheros de configuración:
\begin{itemize}
    \item \textbf{Jugadores:} Catálogo de 36 personajes históricos (4 porteros, 9 defensas, 13 medios y 10 delanteros) con estadísticas base y habilidades propias.
    \item \textbf{Rarezas:} Clasificación en Leyenda, Épico, Raro y Común, determinando el número de habilidades accesibles y la probabilidad de obtención en los contratos.
    \item \textbf{Entrenadores:} Dos figuras técnicas (Cholo Simeone y Luis Aragonés) que otorgan bonificaciones pasivas y disponen de una habilidad definitiva (\textit{ulti}).
    \item \textbf{Armas y Habilidades:} 24 habilidades ofensivas, defensivas y de movilidad con cinco niveles de progresión, limitadas por restricciones temporales de aparición (\texttt{desdeMinuto}) y tiempos de recarga (\texttt{cdMin}).
    \item \textbf{Pasivas y Sinergias:} 14 mejoras pasivas y 24 combinaciones estratégicas (\textit{metas}) que otorgan bonificaciones al acumular determinadas habilidades, desbloqueando a su vez 6 poderes definitivos (\textit{ULTRAS}).
    \item \textbf{Rivales y Jefes:} 16 equipos rivales con comportamientos ofensivos diferenciados y 10 jefes principales dotados de tres fases de combate, escoltas y mecánicas de furia.
    \item \textbf{Mundo Procedural:} Generación infinita de mapas basados en semillas aleatorias, con biomas, obstáculos macizos y colisiones basadas en la posición de los pies de los personajes.
    \item \textbf{Gestión de Club y Servidor:} Persistencia de plantillas, contratos de fichaje adquiridos mediante fichas obtenidas al vencer jefes, y puntos de entrenamiento asignados a cuatro atributos principales (Vida, Defensa, Disparo y Suerte).
    \item \textbf{Logros y Sonido:} 29 logros vinculados a hitos históricos del club que desbloquean equipaciones especiales, acompañados de un motor de audio sintetizado por Web Audio (música chiptune y efectos de grada).
\end{itemize}

\section{Arquitectura del juego}

La arquitectura del juego separa estrictamente la simulación pura y determinista de la capa de presentación y renderizado, tal como esquematiza la Figura \ref{fig:arquitectura-juego}.

\begin{figure}[htbp]
\centering
\begin{tikzpicture}[
    box/.style={rectangle, draw=tema, fill=gris95, rounded corners, text width=5.2cm, align=center, minimum height=0.9cm, font=\small},
    arrow/.style={-{Stealth[length=2.5mm]}, thick, color=tema}
]
    \node[box] (sim) {\textbf{Simulación pura (sin DOM)}\\\texttt{mundo.js}, paso fijo 1/60 s, semilla, armas, jefes};
    \node[box, below=0.6cm of sim] (pres) {\textbf{Presentación y Renderizado}\\\texttt{juego.js}, \texttt{render.js}, sprites 16×20, Web Audio};
    \node[box, below=0.6cm of pres] (srv) {\textbf{Servidor y Anti-trampas}\\\texttt{juegoService.js}, persistencia, validación de partidas};

    \draw[arrow] (sim) -- (pres);
    \draw[arrow] (pres) -- node[right, font=\footnotesize] {POST /api/juego} (srv);
\end{tikzpicture}
\caption{Arquitectura del juego: separación entre simulación, presentación y servidor.}
\label{fig:arquitectura-juego}
\end{figure}

El motor de simulación opera mediante un paso de tiempo fijo de 1/60 s y un generador de números pseudoaleatorios con semilla, garantizando la reproducibilidad exacta de las partidas sin depender de la tasa de fotogramas por segundo (\textit{FPS}) del navegador.

\section{Servidor y validación anti-trampas}

Para preservar la integridad del ranking de la grada y evitar manipulaciones en el progreso del club, el servidor implementa un riguroso mecanismo anti-trampas:
\begin{itemize}
    \item \textbf{Control de partidas en curso:} Al iniciar una partida, el servidor registra el identificador y el reloj de inicio.
    \item \textbf{Validación de resultados:} Al finalizar, se rechazan resultados con duraciones imposibles, tiempos superiores al transcurrido o métricas de jefes, nivel y bajas que excedan los límites biológicos y estadísticos admisibles por el motor.
    \item \textbf{Concurrencia optimista:} Las operaciones de actualización sobre el perfil del club emplean un control de versión numérico (\texttt{version}) para evitar condiciones de carrera en peticiones simultáneas de fichajes o entrenamiento.
\end{itemize}

\section{Logros y evolución por fases}

El sistema de logros registra hitos estadísticos y ofrece contexto histórico sobre el club. El desarrollo del Sprint 2 se estructuró secuencialmente en seis fases operativas en DoneTonic (P20 a P48), abarcando desde el rediseño inicial de la interfaz y el motor del \textit{survivor}, pasando por el club persistente en servidor, la introducción de sinergias y mundos procedurales, hasta el pulido de rarezas, equilibrado de combate, integración de logros y refuerzo de la calidad técnica.
